% =============================================================================
% Q6 -- individual assignment: Abdulrahman Kobeissi
% Parts 6.1--6.3. The final-product vision and schematic are developed later.
% =============================================================================

\providecommand{\qyes}{%
  \colorbox{AUBGreen}{\textcolor{white}{\bfseries\ensuremath{\checkmark}}}}
\providecommand{\qpartial}{%
  \colorbox{AUBYellow}{\textcolor{white}{\bfseries\ensuremath{\sim}}}}
\providecommand{\qgap}{%
  \colorbox{AUBCoral}{\textcolor{white}{\bfseries\ensuremath{\times}}}}

\clearpage
\hypertarget{q6-overview}{}
\sectiontitle{Q6. Individual Assignment -- Abdulrahman Kobeissi}

\qtwosubsection{6.1 Three comparable state-of-the-art technologies}

{\fontsize{9.15}{10.75}\selectfont
The comparison was restricted to systems that solve one of the three central
parts of the present concept: multimodal interface sensing, physical liner
integration, or spatial pressure mapping. The selected technologies are not
presented as identical competitors. Together, they expose which functional
blocks already have credible solutions and which combination remains specific
to the proposed modular 5.2~mm PDMS liner. A commercial screen also considered
the current Tekscan F-Socket system, which uses a paper-thin trimmable sensor,
wired USB acquisition at up to 160~Hz, and real-time 2D/3D pressure software
\cite{tekscanfsocket2026}. The pliance-PMS system was selected for the detailed
commercial comparison because its official documentation discloses sensor
technology, element geometry, range, calibration, channel capacity, telemetry,
and mapping functions \cite{novelpliance2019}.\par}

\vspace{0.45em}
\begin{center}
{\fontsize{7.3}{8.45}\selectfont
\qyes\ \textbf{green check:} function demonstrated in both systems\qquad
\qpartial\ \textbf{yellow partial:} related function, but adaptation or proof is required\qquad
\qgap\ \textbf{red cross:} function absent or deliberately excluded}
\end{center}

\begingroup
\fontsize{7.45}{8.75}\selectfont
\setlength{\tabcolsep}{2.6pt}\renewcommand{\arraystretch}{1.05}
\begin{tabularx}{\linewidth}{>{\raggedright\arraybackslash\bfseries\color{AUBRed}}p{3.0cm}Y Y Y}
\rowcolor{AUBRed}
\color{white}\textbf{Technology} & \color{white}\textbf{Implemented architecture} &
\color{white}\textbf{Evidence already demonstrated} &
\color{white}\textbf{Main boundary relative to our concept} \\
Kwak \textit{et al.}, 2020 & Millimetre-scale 3D Au/PI strain-gauge pressure
sensor plus on-chip temperature sensor; battery-free inner device; paired
NFC/BLE outer module and portable GUI. & Linear pressure response across the
prosthetic range; real-time pressure and temperature at multiple landmarks;
feasibility in transtibial and transfemoral users \cite{q3kwak2020}. & No
humidity measurement, no directional shear output, no liner-wide interpolated
map, and no printed silver-on-PDMS integration. \\
\rowcolor{AUBPale}
Patern\`o \textit{et al.}, 2020 & Twelve commercial Hygrochron iButtons embedded
in CAD-defined blind enclosures in a personalised cryogenically machined
neoprene liner. & Calibrated temperature/RH sensing, anatomical alignment,
wearable liner manufacture, and a 130-min single-participant protocol
\cite{paterno2020}. & Discrete 5.89-mm-thick loggers, offline retrieval,
temperature/RH only, no printed electrodes, and no real-time spatial fusion. \\
novel pliance-PMS & Flexible capacitive pressure pads connected to a wearable
analyser, SD/Bluetooth telemetry, homogeneous-pressure calibration, and 2D/3D
expert software. & Up to 1024 elements; disclosed pads contain 18--90
10$\times$10~mm sensels, are below 1~mm uncoated thickness, cover 10--600~kPa,
and support static/dynamic socket evaluation \cite{novelpliance2019}. & Mature
pressure-only measurement apparatus rather than a multimodal sensing liner;
no printed Ag/PDMS, temperature, RH, or directional shear channel. \\
\end{tabularx}
\endgroup

\vspace{0.55em}
\noindent
\begin{minipage}[t]{0.32\linewidth}\centering
\includegraphics[width=\linewidth,height=0.145\textheight,keepaspectratio]
  {paper_sources/q5_figures/crops/kwak2020_fig1.png}\par
{\fontsize{6.55}{7.45}\selectfont\raggedright
Kwak Fig.~1: sensors, NFC/BLE modules, and the three-block wireless data path
\cite{q3kwak2020}.\par}
\end{minipage}\hfill
\begin{minipage}[t]{0.32\linewidth}\centering
\includegraphics[width=\linewidth,height=0.145\textheight,keepaspectratio]
  {paper_sources/figures_selected/paterno_fig3_liner_design_and_prototype.png}\par
{\fontsize{6.55}{7.45}\selectfont\raggedright
Patern\`o Fig.~3: CAD, machined liner halves, logger enclosures, and final liner
\cite{paterno2020}.\par}
\end{minipage}\hfill
\begin{minipage}[t]{0.32\linewidth}\centering
\includegraphics[width=\linewidth,height=0.145\textheight,keepaspectratio]
  {paper_sources/q6_comparators/crops/pliance_socket_sensors.png}\par
{\fontsize{6.55}{7.45}\selectfont\raggedright
pliance-PMS datasheet: flexible capacitive pads mounted at selected socket
locations \cite{novelpliance2019}.\par}
\end{minipage}

\vfill
{\fontsize{6.75}{7.75}\selectfont\color{AUBGray}
\textbf{Figure Q6-1. Comparator selection.} The three systems were selected by
functional proximity, not visual similarity. Paper figures retain their
original citations; the commercial image is an attributed product-datasheet
excerpt. Detailed evidence appears on the following three pages.\par}

% -----------------------------------------------------------------------------
\clearpage
\hypertarget{q6-kwak}{}
\sectiontitle{Q6.1 Comparator 1 -- Kwak Wireless Multimodal System}

\begin{minipage}[t]{0.36\linewidth}
\centering
\includegraphics[width=\linewidth,height=0.275\textheight,keepaspectratio]
  {paper_sources/q5_figures/crops/kwak2020_fig1.png}\par
{\fontsize{6.65}{7.55}\selectfont\raggedright
\textbf{Published system.} Battery-free pressure/temperature devices sit on
the residual-limb skin; a battery-powered NFC/BLE module sits outside the
socket and relays the data to a portable GUI. Original Fig.~1
\cite{q3kwak2020}.\par}
\end{minipage}\hfill
\begin{minipage}[t]{0.61\linewidth}
{\fontsize{7.6}{8.85}\selectfont
\textbf{Technical architecture.}\par
\begin{itemize}[leftmargin=1.15em,itemsep=0.18em,topsep=0.25em]
\item The pressure transducer is a buckled 3D Au/PI strain-gauge structure,
approximately 1.95$\times$1.95~mm laterally and 0.5~mm high, supported by a
strain-limiting frame rather than printed directly on a liner.
\item The wireless device uses an fPCB, NFC power/data at 13.56~MHz, and a
10-Hz on-chip ADC. The external module forwards data by BLE over more than
10~m; its two-AAA-cell implementation was reported for about 2~h operation.
\item Reported pressure performance included $R^2>0.99$, hysteresis error below
0.02\%, negligible static drift, and a range designed to span interface loads
up to approximately 350~kPa. The design intentionally suppresses shear
sensitivity rather than measuring shear.
\item Multi-site pressure and temperature were demonstrated on two prosthesis
users. The authors explicitly identify humidity, shear, larger participant
samples, and automatic fit evaluation as extensions rather than completed
functions \cite{q3kwak2020}.
\end{itemize}}
\end{minipage}

\vspace{0.55em}
\begingroup
\fontsize{7.15}{8.35}\selectfont
\setlength{\tabcolsep}{2.4pt}\renewcommand{\arraystretch}{1.04}
\begin{tabularx}{\linewidth}{>{\raggedright\arraybackslash\bfseries\color{AUBRed}}p{2.55cm}Y >{\centering\arraybackslash}p{1.55cm}Y}
\rowcolor{AUBRed}
\color{white}\textbf{Feature} & \color{white}\textbf{Kwak implementation} &
\color{white}\textbf{Status} & \color{white}\textbf{Direct comparison with our 5.2 mm slab} \\
Normal pressure & 3D metal strain-gauge transducer with silicone encapsulation
and an independent calibration. & \qyes\par shared & Our pressure block also
requires normal-load calibration, but uses four printed capacitive sectors and
must preserve that response after PDMS molding; see
\reportlink{q2-pressure-shear}{Q2 Section 2.3}. \\
\rowcolor{AUBPale}
Temperature & On-chip temperature channel on the same wireless device. &
\qyes\par shared & Both systems measure temperature, but ours uses a printed
resistive meander that needs slab-specific TCR, self-heating, and bending tests
in \reportlink{q5-microclimate}{Q5 Section 5.4}. \\
Directional shear & Designed to be insensitive to shear; no signed $x/y$
output. & \qgap\par absent & Directional shear is a required independent modality
in our four-capacitor design, making cross-axis separation a major additional
burden in \reportlink{q5-mechanical}{Q5 Section 5.3}. \\
\rowcolor{AUBPale}
Humidity/sweat & Short sweat/high-humidity exposure was used as a robustness
check, not an RH measurement. & \qgap\par absent & Our design includes a printed
IDE humidity sensor and explicitly separates vapour calibration from liquid
sweat protection. \\
Physical integration & Small devices adhere to skin beneath a gel liner; outer
reader modules align across the socket. & \qpartial\par partial & This minimizes
in-socket batteries, but differs from sensors printed on modular PDMS slabs
that become part of the liner structure. \\
\rowcolor{AUBPale}
Acquisition & Real-time, battery-free inner sensor with NFC/BLE relay and GUI. &
\qyes\par reusable & The three-block sensor--gateway--display architecture is a
strong precedent. Our implementation still must prove synchronized acquisition
for four signal types and multiple slabs. \\
Spatial interpretation & Multiple selected anatomical sites; no continuous
liner-wide interpolation or displayed coverage mask. & \qpartial\par partial &
Our coordinate manifest and blinded localization test extend the multinodal
idea into a spatial interface-condition map; see
\reportlink{q5-integration-mapping}{Q5 Section 5.5}. \\
Clinical evidence & Feasibility on two users, without a discomfort-score model
or large statistical study. & \qpartial\par partial & It supports feasibility of
wearable measurements, not validation of our planned discomfort-risk gradient.
\end{tabularx}
\endgroup

\vfill
\sourceanchor{Kwak \textit{et al.} is the closest comparator for wireless
multimodal electronics, but not for printed manufacture, shear, humidity, or
spatial fusion \cite{q3kwak2020}.}

% -----------------------------------------------------------------------------
\clearpage
\hypertarget{q6-paterno}{}
\sectiontitle{Q6.1 Comparator 2 -- Patern\`o Personalised Sensorised Liner}

\begin{minipage}[t]{0.38\linewidth}
\centering
\includegraphics[width=\linewidth,height=0.27\textheight,keepaspectratio]
  {paper_sources/figures_selected/paterno_fig3_liner_design_and_prototype.png}\par
{\fontsize{6.65}{7.55}\selectfont\raggedright
\textbf{Published system.} CAD geometry, cryogenically machined neoprene liner
halves, twelve embedded iButtons, protective fabric, and the final personalised
liner. Original Fig.~3 \cite{paterno2020}.\par}
\end{minipage}\hfill
\begin{minipage}[t]{0.59\linewidth}
{\fontsize{7.55}{8.8}\selectfont
\textbf{Technical architecture.}\par
\begin{itemize}[leftmargin=1.15em,itemsep=0.18em,topsep=0.25em]
\item Twelve DS1923-F5 Hygrochron data loggers measured temperature and RH.
Each stainless-steel logger was 16.25~mm in diameter and 5.89~mm thick and was
placed in a CAD-defined enclosure for direct skin contact.
\item The non-waterproof loggers were covered by breathable waterproof fabric.
Each was chamber-calibrated over 25--45~$^\circ$C and 20--95\%RH before use.
\item Residual-limb and liner scans were landmark-aligned; the liner was split
into two parts and cryogenically CNC-machined from closed-cell neoprene. The
reported build took about 4~h and weighed 107.3~g.
\item One transtibial participant completed 50~min rest, 30~min walking, and
50~min post-activity rest. Data were logged locally and downloaded through a
1-Wire USB interface rather than streamed as a real-time map
\cite{paterno2020}.
\end{itemize}}
\end{minipage}

\vspace{0.55em}
\begingroup
\fontsize{7.15}{8.35}\selectfont
\setlength{\tabcolsep}{2.4pt}\renewcommand{\arraystretch}{1.04}
\begin{tabularx}{\linewidth}{>{\raggedright\arraybackslash\bfseries\color{AUBRed}}p{2.55cm}Y >{\centering\arraybackslash}p{1.55cm}Y}
\rowcolor{AUBRed}
\color{white}\textbf{Feature} & \color{white}\textbf{Patern\`o implementation} &
\color{white}\textbf{Status} & \color{white}\textbf{Direct comparison with our 5.2 mm slab} \\
True liner integration & Sensor pockets were created in a personalised liner
derived from patient scans. & \qyes\par shared & This is the closest precedent for
preserving fit while embedding sensors. Our route uses repeated modular
silicone slabs rather than a split CNC-neoprene liner; see
\reportlink{q2-whole-liner}{Q2 Section 2.7}. \\
\rowcolor{AUBPale}
Temperature and RH & Twelve calibrated dual-function loggers at registered
locations. & \qyes\par shared & These are two of our four modalities and establish
the need for chamber calibration and location identity. \\
Pressure and shear & Not instrumented in the reported liner. & \qgap\par absent &
Our pressure/directional-shear pair adds load transfer, cross-axis separation,
and mechanical hysteresis risks that this system did not need to solve. \\
\rowcolor{AUBPale}
Sensor form factor & Rigid 5.89-mm-thick commercial logger in a blind enclosure.
& \qpartial\par partial & The thickness is close to our complete 5.2-mm slab, but
our printed sensor is distributed within the slab instead of occupying a rigid
local cavity. \\
Moisture protection & Breathable waterproof fabric protected the logger while
allowing environmental exchange. & \qpartial\par reusable & This is a functional
precedent for a vapour-access/liquid-exclusion barrier, but H1/H2 must verify it
for the printed IDE and exposed silver interfaces. \\
\rowcolor{AUBPale}
Manufacture & Scan alignment, CAD subtraction, split design, and cryogenic CNC
machining of neoprene. & \qpartial\par partial & The digital personalization
workflow is reusable; the material, printing, cure, electrical routing, and
slab-to-liner assembly are different. \\
Acquisition and mapping & Local logging and post-test download; graphs retain
sensor location but are not a live surface reconstruction. & \qgap\par absent &
Our concept requires synchronized acquisition, a CAD coordinate manifest, and
a real-time map with visible measured sites and uncertainty. \\
\rowcolor{AUBPale}
Human evidence & One participant and one liner architecture; comfort and
thermal outcomes were reported. & \qpartial\par partial & Valuable wearable-use
evidence, but insufficient to validate universal thresholds or our combined
discomfort-risk score. \\
\end{tabularx}
\endgroup

\vfill
\sourceanchor{Patern\`o \textit{et al.} provides the strongest liner-integration
and microclimate precedent. Its published limitations closely define what our
pressure, shear, printed manufacture, and real-time mapping blocks add
\cite{paterno2020}.}

% -----------------------------------------------------------------------------
\clearpage
\hypertarget{q6-pliance}{}
\sectiontitle{Q6.1 Comparator 3 -- novel pliance-PMS Prosthesis System}

\begin{minipage}[t]{0.43\linewidth}
\centering
\includegraphics[width=\linewidth,height=0.16\textheight,keepaspectratio]
  {paper_sources/q6_comparators/crops/pliance_socket_sensors.png}\par
\vspace{0.3em}
\includegraphics[width=\linewidth,height=0.15\textheight,keepaspectratio]
  {paper_sources/q6_comparators/crops/pliance_software_and_map.png}\par
{\fontsize{6.55}{7.45}\selectfont\raggedright
\textbf{Commercial system.} Flexible socket-mounted sensor pads, acquisition
software, and registered 3D pressure display from the official pliance-PMS
prosthesis datasheet \cite{novelpliance2019}.\par}
\end{minipage}\hfill
\begin{minipage}[t]{0.54\linewidth}
{\fontsize{7.5}{8.75}\selectfont
\textbf{Disclosed technical architecture.}\par
\begin{itemize}[leftmargin=1.15em,itemsep=0.17em,topsep=0.25em]
\item Flexible capacitive pads are placed at selected anatomical regions and
connected to a pliance-PMS analyser. The analyser can scan up to 1024 elements
and transmit by Bluetooth or store data on an SD card.
\item Listed pad examples contain 18, 45, or 90 sensels. Each sensel is
10$\times$10~mm; the uncoated pad thickness is below 1~mm, elasticity is about
4\%, and the disclosed pressure range is 10--600~kPa.
\item The product literature reports acquisition up to 20,000 sensor samples
per second, selectable matrix localization, 2D/3D display, and synchronization
with video and EMG.
\item A trublu calibration device applies homogeneous air pressure to all
elements for simultaneous calibration. The 2019 datasheet also states that the
product was not certified as a European medical product at that time
\cite{novelpliance2019}.
\end{itemize}}
\end{minipage}

\vspace{0.5em}
\begingroup
\fontsize{7.1}{8.3}\selectfont
\setlength{\tabcolsep}{2.35pt}\renewcommand{\arraystretch}{1.035}
\begin{tabularx}{\linewidth}{>{\raggedright\arraybackslash\bfseries\color{AUBRed}}p{2.55cm}Y >{\centering\arraybackslash}p{1.55cm}Y}
\rowcolor{AUBRed}
\color{white}\textbf{Feature} & \color{white}\textbf{pliance-PMS implementation} &
\color{white}\textbf{Status} & \color{white}\textbf{Direct comparison with our 5.2 mm slab} \\
Pressure range & Calibrated capacitive pads covering 10--600~kPa. &
\qyes\par shared & The commercial range contains our 0--250~kPa Q4 target, but
our four-sector printed cell needs its own calibration, saturation, and
hysteresis evidence. \\
\rowcolor{AUBPale}
Spatial density & 10$\times$10~mm sensels; selected pads provide 18--90
elements and the analyser supports up to 1024. & \qyes\par benchmark & This is a
stronger pressure-density benchmark than the current sparse multimodal slab
layout and exposes potential missed-hotspot risk. \\
Directional shear & No signed tangential-force channel in the disclosed PMS
sensor set. & \qgap\par absent & Our design aims to recover $x/y$ shear from
opposed capacitor pairs while maintaining normal-pressure separation. \\
\rowcolor{AUBPale}
Temperature and RH & Not measured by the PMS pads. & \qgap\par absent & Our liner
adds separate printed meander and IDE modalities, increasing clinical context
but also calibration and synchronization complexity. \\
Geometry and comfort & Thin, flexible, elastic pads conform to curved sites but
remain added measurement hardware at selected locations. & \qpartial\par partial &
Our slab is thicker but intended to become part of the liner structure, with
sensor placement fixed by CAD rather than temporary pad installation. \\
\rowcolor{AUBPale}
Calibration & Homogeneous pneumatic calibration across all capacitive elements.
& \qyes\par reusable & The independent pressure reference and simultaneous channel
release are directly reusable concepts for M1, although the printed stack and
force-transfer geometry differ. \\
Mapping software & Matrix localization, 2D/3D pressure display, Bluetooth, and
video/EMG synchronization. & \qyes\par benchmark & This is the closest mature map
benchmark. Our map must add four physical-unit layers, signed shear, confidence
masks, and calibration metadata without averaging conflicting modalities. \\
\rowcolor{AUBPale}
Clinical interpretation & Quantifies interface pressure for socket evaluation;
the datasheet does not validate a multimodal discomfort score. &
\qpartial\par partial & Our display must remain an interface-condition map until
participant-reported discomfort and clinically justified thresholds are added.
\end{tabularx}
\endgroup

\vspace{0.45em}
\evidencebox{Comparison verdict:}{Each comparator already solves a substantial
functional block. The defensible project distinction is the proposed
combination of printed pressure, directional shear, temperature, and humidity
on coordinate-registered 5.2~mm PDMS slabs with modality-preserving software.
This is an engineering differentiation statement, not a patentability or
clinical-superiority claim.}

% -----------------------------------------------------------------------------
\clearpage
\hypertarget{q6-challenges-materials}{}
\sectiontitle{Q6.2 Challenges -- Materials, Printing, and Sensor Physics}

{\fontsize{8.8}{10.35}\selectfont
\textbf{Evidence boundary.} Two fabrication milestones are confirmed: the
sensor artwork was printed successfully on PET as the accessible first-stage
process check, and a 5.2~mm PDMS slab was molded. These achievements do not yet
establish direct-on-PDMS print quality, a calibrated multimodal slab, or
whole-liner durability. The challenges below are therefore stated as
engineering questions with published causes, project-specific consequences,
and linked validation routes---not as claims that the prototype has failed.\par}

\vspace{0.42em}
\begingroup
\fontsize{7.15}{8.35}\selectfont
\setlength{\tabcolsep}{2.25pt}\renewcommand{\arraystretch}{1.045}
\begin{tabularx}{\linewidth}{>{\raggedright\arraybackslash\bfseries\color{AUBRed}}p{2.38cm}Y Y Y}
\rowcolor{AUBRed}
\color{white}\textbf{Challenge} &
\color{white}\textbf{What the literature establishes} &
\color{white}\textbf{Present design implication} &
\color{white}\textbf{Linked evidence route} \\
PDMS wetting and solvent swelling &
Oxygen-plasma treatment enabled continuous AgNP printing on PDMS, while the
same printer settings produced an apparent profile of about 3.8~$\mu$m on
PDMS versus 0.6~$\mu$m on PET; much of the difference was attributed to
solvent-induced PDMS swelling \cite{albrecht2018stretch}. &
PET confirms artwork transfer and printer operation only. It cannot predict
trace height, spreading, resistance, or sensor calibration on molded PDMS. &
Measure contact angle, line width, dry/equilibrated profile, continuity, and
sheet resistance using the process controls in
\reportlink{q2-pdms-printing}{Q2 Section 2.5} and release gate
\reportlink{q5-release}{Q5 Section 5.2}. \\
\rowcolor{AUBPale}
Silver adhesion, cracking, and stretch &
Direct PDMS routes use surface activation, wider trace geometry,
pre-stretching, or semi-embedded conductors to improve continuity under
deformation \cite{sun2016,q4abukhalaf2018,q4abukhalaf2019}. &
Successful printing is not equivalent to a mechanically stable electrode.
Compression, flexion, shear, cable pull, and cure history can change resistance
or create intermittent opens. &
Use witness traces and full sensor coupons; record resistance during bending
and cyclic loading, adhesion, crack images, and recovery as linked in
\reportlink{q4-fabrication}{Q4 Section 4.3} and
\reportlink{q5-durability}{Q5 Section 5.6}. \\
5.2~mm force-transfer stack &
The four-capacitor shear architecture obtains direction from differential
capacitance; dielectric thickness, overlap, spacer geometry, normal preload,
and viscoelastic hysteresis all affect the recovered vector
\cite{albrecht2019}. &
The molded thickness is a geometric fact, not evidence that local pressure and
signed shear reach the printed electrodes without attenuation or cross-axis
coupling. &
Calibrate normal load and signed $x/y$ shear at several preloads, then quantify
sensitivity, nonlinearity, hysteresis, drift, and axis coupling in
\reportlink{q5-mechanical}{Q5 Section 5.3}. \\
\rowcolor{AUBPale}
Temperature-channel stability &
Printed Ag temperature meanders require trace-specific temperature
coefficient, response-time, bending, humidity, and self-heating
characterization \cite{liew2020,jager2022}. &
The meander may respond to temperature, strain, lead resistance, and changing
thermal contact simultaneously after encapsulation. &
Fit resistance--temperature calibration on the final stack; repeat it dry,
humid, bent, and after cycling using
\reportlink{q2-temperature}{Q2 Section 2.4} and
\reportlink{q5-microclimate}{Q5 Section 5.4}. \\
Humidity: vapour versus liquid sweat &
Printed IDE sensors can measure RH, while embedded-liner work shows that a
breathable waterproof barrier and chamber calibration are necessary for
interpretable skin-interface microclimate data
\cite{aitmammar2020,paterno2020}. &
An open IDE may saturate or leak under liquid sweat; an aggressive barrier may
prevent liquid failure but slow or bias the vapour response. &
Run controlled RH steps separately from metered artificial sweat; report
response/recovery time, hysteresis, leakage, and post-drying recovery in
\reportlink{q2-humidity}{Q2 Section 2.4} and
\reportlink{q5-microclimate}{Q5 Section 5.4}. \\
\rowcolor{AUBPale}
Hydrogel as a later interface variable &
Hydrogel--PDMS bonding and sweat uptake are feasible, but water uptake depends
on thickness and time and changes mechanics and transport
\cite{yuk2016,lu2025}. No reviewed study validates a hydrogel over this exact
four-modality prosthetic slab. &
A hydrogel may soften contact, but it may also store sweat, delay RH and
temperature, redistribute pressure, alter friction/shear, or expose silver to
water and ions. It is not the baseline architecture. &
First validate bare PDMS H0; then compare H1 islands and the H2 windowed sheet
against H0. The structure and rejection criteria are in
\reportlink{q2-hydrogel-interactions}{Q2 Section 2.6.1} and
\reportlink{q2-hydrogel-structures}{Q2 Section 2.6.2}. \\
\end{tabularx}
\endgroup

\vspace{0.4em}
\begin{center}
\begin{tikzpicture}[>=stealth,node distance=3.7mm]
  \tikzstyle{qstage}=[draw=AUBLine,rounded corners=1.2pt,minimum height=7.5mm,
    text width=28mm,align=center,font=\fontsize{6.6}{7.4}\selectfont]
  \node[qstage,fill=AUBGreen!14] (pet) {confirmed\\PET print};
  \node[qstage,fill=AUBGreen!14,right=of pet] (mold) {confirmed\\5.2 mm PDMS mold};
  \node[qstage,fill=AUBYellow!18,right=of mold] (print) {open evidence\\direct PDMS print};
  \node[qstage,fill=AUBYellow!18,right=of print] (cal) {open evidence\\multimodal calibration};
  \node[qstage,fill=AUBCoral!13,right=of cal] (liner) {later gate\\whole-liner survival};
  \draw[->,thick,color=AUBRed] (pet)--(mold);
  \draw[->,thick,color=AUBRed] (mold)--(print);
  \draw[->,thick,color=AUBRed] (print)--(cal);
  \draw[->,thick,color=AUBRed] (cal)--(liner);
\end{tikzpicture}
\end{center}
{\fontsize{6.9}{7.9}\selectfont\color{AUBGray}
\textbf{Figure Q6-2. Evidence progression for the material stack.} Green marks
documented project stages; yellow marks the next measurement gates; the final
liner gate follows only after coupon and slab evidence. This sequence is
consistent with the fabrication evidence reviewed in
\reportlink{q2-pdms-printing}{Q2 Section 2.5}.\par}

% -----------------------------------------------------------------------------
\clearpage
\hypertarget{q6-challenges-system}{}
\sectiontitle{Q6.2 Challenges -- Integration, Mapping, Durability, and Use}

\begingroup
\fontsize{7.22}{8.43}\selectfont
\setlength{\tabcolsep}{2.2pt}\renewcommand{\arraystretch}{1.05}
\begin{tabularx}{\linewidth}{>{\raggedright\arraybackslash\bfseries\color{AUBRed}}p{2.35cm}Y Y Y}
\rowcolor{AUBRed}
\color{white}\textbf{Challenge} &
\color{white}\textbf{Evidence from comparable systems} &
\color{white}\textbf{Specific risk in this design} &
\color{white}\textbf{Required control and report link} \\
Sensor identity and CAD coordinates &
Patern\`o used scan-aligned anatomical landmarks and CAD-defined pockets;
Kwak retained the identity of multiple selected sites
\cite{paterno2020,q3kwak2020}. &
A correct value attached to the wrong slab, modality, or coordinate creates a
convincing but false hotspot. Curvature and assembly can shift the printed site
away from its nominal CAD point. &
Create a permanent channel--slab--modality--coordinate manifest and perform
blinded stimulus localization in
\reportlink{q5-integration-mapping}{Q5 Section 5.5}. \\
\rowcolor{AUBPale}
Sparse coverage and missed hotspots &
Commercial pliance-PMS arrays disclose 10$\times$10~mm sensels and up to 1024
elements; Turner used 144 pressure sensors arranged as twelve strips
\cite{novelpliance2019,q3turner2021}. &
The proposed multimodal slabs are intentionally sparser. A hotspot between
sensors may be attenuated, shifted, or entirely absent, especially near slab
boundaries. &
Test centres, between-sensor points, and boundaries; report nearest-neighbour
pitch, coverage mask, median/95th-percentile localization error, and density
trade-offs in \reportlink{q3-metrics}{Q3 Section 3.7} and
\reportlink{q5-integration-mapping}{Q5 Section 5.5}. \\
Interpolation and false visual confidence &
Turner demonstrated real-time 3D socket visualization using radial-basis
interpolation, proving the display architecture but not that interpolated
pixels are measured data \cite{q3turner2021}. &
A smooth colour gradient can hide sparse coverage and imply certainty between
physical sensors. Combining different units into one colour also risks
misleading averaging. &
Display pressure, signed shear, temperature, and RH separately first; retain
sensor markers, unsupported-region masks, method/version metadata, and
confidence flags. See \reportlink{q3-measurement-models}{Q3 Section 3.5} and
\reportlink{q4-mapping}{Q4 Section 4.6}. \\
\rowcolor{AUBPale}
Sampling and multimodal synchronization &
Kwak used real-time multimodal acquisition, Turner sampled pressure at
200~Hz, while Patern\`o's microclimate loggers were downloaded after the
activity protocol \cite{q3kwak2020,q3turner2021,paterno2020}. &
Mechanical events are fast; temperature and RH are slower. Unsynchronized
streams can incorrectly associate a pressure/slip event with a later
microclimate response. &
Use one monotonic clock; preserve raw timestamps and modality-specific rates;
define resampling windows explicitly and verify combined-condition timing in
\reportlink{q5-integration-mapping}{Q5 Section 5.5}. \\
Cross-sensitivity and compensation &
Published sensors are normally calibrated by modality. The Kwak architecture
suppressed shear sensitivity; the printed shear cell relies on differential
responses; humidity and temperature outputs depend on their environments
\cite{q3kwak2020,albrecht2019,aitmammar2020,jager2022}. &
Pressure, shear, bending, temperature, moisture, cable motion, and hydrogel
swelling may change multiple electrical channels at once. A compensation
model could overfit and remove real events. &
Apply one controlled input at a time before combined tests; fit compensation
on training runs and verify it on held-out combinations. Preserve raw and
corrected values as specified in
\reportlink{q5-integration-mapping}{Q5 Section 5.5}. \\
\rowcolor{AUBPale}
Cyclic drift and whole-liner routing &
Printed Ag/PDMS conductors can survive demanding stretch tests under published
processes, but performance depends on ink, surface treatment, geometry, and
loading mode \cite{albrecht2018stretch,sun2016,q4abukhalaf2018}. &
Repeated compression, bending, donning, sweat, cure interfaces, connectors, and
strain relief may change baselines even when an isolated trace remains
continuous. &
Repeat calibration after staged mechanical and wet/dry exposure; inspect
cracks and adhesion; quantify zero shift, span change, hysteresis, and channel
loss in \reportlink{q5-durability}{Q5 Section 5.6} before the
\reportlink{q2-whole-liner}{Q2 Section 2.7} integration step. \\
Phantom evidence and clinical meaning &
Published systems support interface measurement and controlled socket
evaluation, yet the human evidence in close comparators is small and a recent
systematic review emphasizes limitations in using pressure alone for clinical
decisions \cite{q3kwak2020,paterno2020,q3armitage2026}. &
A phantom can verify loading, moisture, temperature, localization, and
repeatability; it cannot prove perceived discomfort or tissue-injury risk.
Thresholds based only on engineering judgment must not be presented as
clinical truth. &
Use the modular phantom protocols in
\reportlink{q5-phantom}{Q5 Section 5.7}. Call the output an
\emph{interface-condition map}; reserve \emph{discomfort-risk} for later
participant-reported and clinically reviewed validation. \\
\end{tabularx}
\endgroup

\vspace{0.45em}
\begin{center}
\begin{tikzpicture}[x=1cm,y=1cm,>=stealth]
  \node[draw=AUBLine,fill=AUBPale,rounded corners,text width=31mm,align=center,
    minimum height=9mm,font=\fontsize{6.8}{7.7}\selectfont] (raw) at (0,0)
    {raw channel\\value + unit + time};
  \node[draw=AUBLine,fill=AUBPale,rounded corners,text width=31mm,align=center,
    minimum height=9mm,font=\fontsize{6.8}{7.7}\selectfont] (manifest) at (4.1,0)
    {calibration version\\slab + CAD coordinate};
  \node[draw=AUBLine,fill=AUBYellow!15,rounded corners,text width=31mm,align=center,
    minimum height=9mm,font=\fontsize{6.8}{7.7}\selectfont] (map) at (8.2,0)
    {physical-unit map\\markers + coverage mask};
  \node[draw=AUBLine,fill=AUBCoral!12,rounded corners,text width=31mm,align=center,
    minimum height=9mm,font=\fontsize{6.8}{7.7}\selectfont] (claim) at (12.3,0)
    {interface-condition\\interpretation + confidence};
  \draw[->,thick,color=AUBRed] (raw)--(manifest);
  \draw[->,thick,color=AUBRed] (manifest)--(map);
  \draw[->,thick,color=AUBRed] (map)--(claim);
  \node[font=\fontsize{6.4}{7.2}\selectfont,text=AUBGray,align=center] at (6.15,-0.85)
    {every transformation remains traceable; interpolation never replaces raw evidence};
\end{tikzpicture}
\end{center}
{\fontsize{6.9}{7.9}\selectfont\color{AUBGray}
\textbf{Figure Q6-3. System evidence chain.} The software can only make a
defensible spatial statement if channel identity, calibration, time, and CAD
position survive every processing stage. The architecture synthesizes the
published acquisition and mapping precedents
\cite{q3kwak2020,paterno2020,q3turner2021}.\par}

\vspace{0.38em}
\evidencebox{Priority conclusion:}{The highest-risk path is not any single
sensor principle. It is preserving calibrated, synchronized, correctly located
four-modality information after direct PDMS printing, 5.2~mm mechanical
integration, repeated loading, moisture exposure, and spatial interpolation.
That path is tested progressively in Q5 rather than deferred to a complete
liner.}

% -----------------------------------------------------------------------------
\clearpage
\hypertarget{q6-team}{}
\sectiontitle{Q6.3 Team Expertise, Ownership, and Remaining Gaps}

{\fontsize{8.9}{10.45}\selectfont
The project already spans biomedical, chemical/material, and mechanical/
electronic work. The table records who owns each functional contribution; it
does not imply that one member works in isolation. Interfaces between roles are
explicit because printing, material selection, mechanics, circuitry, and
software all affect the same calibrated signal.\par}

\vspace{0.42em}
\begingroup
\fontsize{7.55}{8.85}\selectfont
\setlength{\tabcolsep}{2.5pt}\renewcommand{\arraystretch}{1.08}
\begin{tabularx}{\linewidth}{>{\raggedright\arraybackslash\bfseries\color{AUBRed}}p{2.55cm}p{2.6cm}Y Y}
\rowcolor{AUBRed}
\color{white}\textbf{Member} & \color{white}\textbf{Discipline} &
\color{white}\textbf{Primary responsibility} &
\color{white}\textbf{Concrete project outputs and interfaces} \\
Abdulrahman Kobeissi & \shortstack[l]{Biomedical\\Engineering} &
System architecture; all printed-sensor CAD and print work; sensor proportions
and placement; 5.2~mm silicone slab and circuit-layout design; software
algorithms for calibration, interpolation, coordinate mapping, and
interface-condition visualization. &
PET print artwork and execution; CAD sensor/slab designs; multimodal placement
logic; channel-to-coordinate manifest; processing and map pipeline. Interfaces
with Hawraa on PDMS/hydrogel compatibility and Amir on electrical signals,
sampling, and connector definitions. \\
\rowcolor{AUBPale}
Hawraa & Chemistry &
Hydrogel and silicone material science; surface chemistry, bonding, moisture
uptake, protection, and material compatibility. &
PDMS and hydrogel formulation/process input; swelling, absorption, drying,
bonding, and wet/dry compatibility criteria. Interfaces with Abdulrahman on
sensor exposure/window geometry and with Amir on insulation, leakage, and
material effects on the readout. \\
Amir & \shortstack[l]{Mechanical\\Engineering} &
Physical sensor acquisition, mechanical test implementation, signal/data
acquisition, and circuitry/readout design. &
Loading and acquisition hardware; circuit topology; channel excitation and
readout; sampling and synchronization implementation; force/application
fixtures and connector/strain-relief input. Interfaces with Abdulrahman on the
data schema and with Hawraa on encapsulation and wet-operation constraints. \\
\end{tabularx}
\endgroup

\vspace{0.65em}
\qtwosubsection{Expertise not yet represented by a named specialist}

\begingroup
\fontsize{7.3}{8.55}\selectfont
\setlength{\tabcolsep}{2.35pt}\renewcommand{\arraystretch}{1.065}
\begin{tabularx}{\linewidth}{>{\raggedright\arraybackslash\bfseries\color{AUBRed}}p{2.75cm}Y Y}
\rowcolor{AUBRed}
\color{white}\textbf{Missing specialist perspective} &
\color{white}\textbf{Why it matters} &
\color{white}\textbf{How the team has compensated at this stage} \\
Certified prosthetist/orthotist and rehabilitation clinician &
Socket fit, bony landmarks, safe loading, donning constraints, clinical
workflow, and interpretation of whether a measured condition is meaningful to
a prosthesis user require clinical and prosthetic practice. &
Placement and validation are grounded in prosthetic-interface literature and
published liner systems \cite{paterno2020,q3kwak2020,q3tang2022}. Work remains
on benches and phantoms; no human-safety, comfort, or clinical-risk claim is
made. External clinical review is required before participant testing. \\
\rowcolor{AUBPale}
Printed-electronics process/metrology specialist &
AgNP wetting, plasma activation, sintering, profilometry, adhesion, swelling,
crack morphology, and production yield on PDMS are process-specific. PET
success cannot certify them. &
The team separated the PET print stage from direct-PDMS validation, adopted
published process controls \cite{albrecht2018stretch,sun2016,
q4abukhalaf2018,q4abukhalaf2019}, and specified coupons, witness traces,
profilometry/microscopy, resistance tracking, and build records in
\reportlink{q5-release}{Q5 Section 5.2}. \\
Human-factors, biostatistics, and regulatory/quality expertise &
Discomfort thresholds, sample size, uncertainty, usability, traceability,
participant protection, and the boundary between an engineering tool and a
medical claim cannot be settled by sensor calibration alone. &
The report predeclares test gates, uses held-out validation, retains raw and
processed data with calibration versions, distinguishes pass/partial/fail/not
run, and limits current language to an interface-condition map. Participant
and clinical validation remains a separate reviewed stage; see
\reportlink{q4-design-gate}{Q4 Section 4.7} and
\reportlink{q5-reporting}{Q5 Section 5.8}. \\
\end{tabularx}
\endgroup

\vspace{0.55em}
\begin{center}
\begin{tikzpicture}[>=stealth]
  \node[draw=AUBRed,fill=AUBPale,rounded corners,text width=34mm,align=center,
    minimum height=12mm,font=\fontsize{7}{8}\selectfont] (ak) at (0,0)
    {\textbf{Abdulrahman}\\sensor CAD + printing\\placement + algorithms};
  \node[draw=AUBRed,fill=AUBPale,rounded corners,text width=34mm,align=center,
    minimum height=12mm,font=\fontsize{7}{8}\selectfont] (ha) at (5.1,0)
    {\textbf{Hawraa}\\PDMS + hydrogel\\chemistry + compatibility};
  \node[draw=AUBRed,fill=AUBPale,rounded corners,text width=34mm,align=center,
    minimum height=12mm,font=\fontsize{7}{8}\selectfont] (am) at (10.2,0)
    {\textbf{Amir}\\mechanics + circuitry\\signal/data acquisition};
  \node[draw=AUBYellow!80!black,fill=AUBYellow!13,rounded corners,text width=75mm,
    align=center,minimum height=9mm,font=\fontsize{6.8}{7.7}\selectfont] (ext)
    at (5.1,-2.0) {external review gates: prosthetics/clinical practice
    \quad printed-electronics metrology \quad human factors/statistics/quality};
  \draw[<->,thick,color=AUBRed] (ak)--(ha);
  \draw[<->,thick,color=AUBRed] (ha)--(am);
  \draw[<->,thick,color=AUBRed] (ak) to[bend right=17] (am);
  \draw[->,thick,color=AUBYellow!80!black] (ak)--(ext);
  \draw[->,thick,color=AUBYellow!80!black] (ha)--(ext);
  \draw[->,thick,color=AUBYellow!80!black] (am)--(ext);
\end{tikzpicture}
\end{center}
{\fontsize{6.9}{7.9}\selectfont\color{AUBGray}
\textbf{Figure Q6-4. Responsibility and review structure.} Internal ownership
covers the principal engineering blocks; the lower review gate identifies
specialist input required before human testing or clinical interpretation.\par}

\vspace{0.38em}
\evidencebox{Team conclusion:}{The principal prototype can be developed by the
current multidisciplinary team. Literature-based controls and modular bench
tests reduce the consequences of missing specialist expertise, but they do not
replace external prosthetic/clinical, process-metrology, and human-validation
review at the relevant design gates.}

% -----------------------------------------------------------------------------
\clearpage
\hypertarget{q6-integrated-liner-render}{}
\sectiontitle{Q6.4 Integrated Liner Architecture}

{\fontsize{9}{10.6}\selectfont
The consolidated design places the compact multimodal slabs exclusively on the
inner, skin-contacting surface of the liner. The rolled-open view exposes the
distributed sensing array while preserving a smooth silicone exterior. Signal
routing remains enclosed within the liner wall, and the burgundy electronics
enclosure is mounted externally on the lateral side.\par}

\vspace{0.35em}
\begin{center}
\includegraphics[width=\linewidth,height=0.64\textheight,keepaspectratio]
  {assets/q6/smart_liner_inner_sensor_array_rolled_cutaway_v5.png}\par
\end{center}
{\fontsize{7.3}{8.5}\selectfont\color{AUBGray}
\textbf{Figure Q6-5. Integrated smart-liner concept.} The main rolled-open view
shows the dense slab array on the skin-contacting interior and the plain
silicone exterior with its side-mounted electronics enclosure. The exploded
view identifies the layered slab construction, while the distal inset shows
how the inner modules continue toward the closed end of the liner. This is a
project design visualization rather than evidence of a completed whole-liner
prototype.\par}

\vspace{0.45em}
\evidencebox{Architecture represented:}{Inner multimodal sensing coverage;
flexible spacing between complete slabs; protected in-wall signal routing;
plain outer silicone; and an externally accessible lateral electronics pod.}
